% Build Your Own Mini Docker — Paradox 2026 workshop slides.
% Beamer + metropolis, compiled with XeLaTeX (see Makefile).
% Slides are an AID — they hold the headline, you speak the rest. Keep them light.
\documentclass[10pt,aspectratio=169]{beamer}
\usetheme{metropolis}
\metroset{numbering=fraction, progressbar=frametitle, sectionpage=progressbar}

\usepackage{booktabs}
\usepackage{listings}

% Light monospace style for the handful of command snippets we show.
\lstset{
  basicstyle=\ttfamily\small,
  columns=fullflexible,
  keepspaces=true,
  breaklines=true,
  aboveskip=0.4em, belowskip=0.4em,
}
\newcommand{\cmd}[1]{\texttt{\small #1}}

\title{Build Your Own Mini Docker}
\subtitle{Understanding containers from scratch}
\author{Sayan}
\date{Paradox 2026 \quad\textbar\quad IIT Madras \quad\textbar\quad 11 June}

\begin{document}

\maketitle

% ============================================================ HOOK
\begin{frame}[standout]
  What is a container,\\[0.3em]
  \emph{in syscalls?}
\end{frame}

\begin{frame}[fragile]{Before we touch any of that}
  \begin{itemize}
    \item Entry ticket: a \alert{green pre-flight} on your laptop.
    \item Get the code:
  \end{itemize}
  \begin{lstlisting}
git clone https://github.com/sayan-iitm/minidocker
cd minidocker
  \end{lstlisting}
  \begin{itemize}
    \item Fell behind at any point? \cmd{bash demo/cpN/run.sh}.
  \end{itemize}
\end{frame}

% ============================================================ NAMESPACES
\section{Namespaces}

\begin{frame}{One idea, seven times}
  \begin{center}
    \Large A namespace is an \alert{isolated view} of one global thing.
  \end{center}
  \vspace{1em}
  \begin{center}
  \begin{tabular}{ll}
    \toprule
    \textbf{mount}  & the filesystem tree \\
    \textbf{UTS}    & the hostname \\
    \textbf{PID}    & process IDs \\
    \textbf{net}    & interfaces, routes, ports \\
    \textbf{user}   & uid / gid — root inside $\neq$ root outside \\
    \textbf{IPC}    & shared memory, message queues \\
    \textbf{cgroup} & the cgroup tree you can see \\
    \bottomrule
  \end{tabular}
  \end{center}
\end{frame}

\begin{frame}[fragile]{One command runs the lot}
  \begin{lstlisting}
unshare --uts --pid --mount --user --net --ipc --cgroup ...
  \end{lstlisting}
  \begin{itemize}
    \item \cmd{lsns} — the namespaces you're already in.
    \item The next \alert{$\sim$100 lines of bash} just glue these together.
  \end{itemize}
\end{frame}

% ============================================================ THE BUILD
\section{100 lines of bash}

\begin{frame}{Where we're going}
  \begin{enumerate}
    \item \textbf{CP1} — hostname \hfill {\small UTS}
    \item \textbf{CP2} — our own process list \hfill {\small PID + mount}
    \item \textbf{CP3} — drop \cmd{sudo} \hfill {\small user}
    \item \textbf{CP4} — our own root filesystem \hfill {\small pivot\_root}
    \item \textbf{CP5} — a limit the kernel enforces \hfill {\small cgroups}
  \end{enumerate}
  \vspace{1em}
  \begin{center}
    \small then tie it together $\rightarrow$ \cmd{mini-docker.sh}
  \end{center}
  \vfill
  {\footnotesize Lost? \cmd{bash demo/cpN/run.sh} catches you up.}
\end{frame}

% ============================================================ PID 1
\begin{frame}{Ctrl-C did nothing. Why?}
  \begin{itemize}
    \item Inside the container, our shell is \alert{PID 1}.
    \item The kernel \alert{drops} \cmd{SIGINT}/\cmd{SIGTERM} to PID 1 unless it
          installs a handler.
    \item Real runtimes ship a tiny init (\cmd{tini}) that becomes PID 1,
          forwards signals, reaps zombies.
  \end{itemize}
  \vspace{1em}
  \begin{center}
    \Large That is what \alert{\texttt{docker run -{}-init}} is for.
  \end{center}
\end{frame}

% ============================================================ WHAT DOCKER ADDS
\section{What Docker adds on top}

\begin{frame}{Images — OCI}
  \begin{itemize}
    \item A tarball of \alert{layers} + metadata (what to run, env, ports).
    \item Our \cmd{rootfs/} was step one of this, by hand.
  \end{itemize}
\end{frame}

\begin{frame}{Layers — OverlayFS}
  \begin{itemize}
    \item Stack read-only layers; add one thin \alert{writable} layer on top.
    \item Share the base across containers $\Rightarrow$ fast, cheap.
  \end{itemize}
\end{frame}

\begin{frame}{Dockerfile — the build}
  \begin{itemize}
    \item A recipe. Each instruction $\rightarrow$ one cached layer.
    \item Reproducible images instead of "works on my machine".
  \end{itemize}
\end{frame}

\begin{frame}{Registries}
  \begin{itemize}
    \item \cmd{push} / \cmd{pull} images by name + tag.
    \item Docker Hub, GHCR, your own — just content-addressed blobs.
  \end{itemize}
\end{frame}

\begin{frame}{Networking \& security}
  \begin{itemize}
    \item bridge + \cmd{veth} pairs + \cmd{iptables} $\rightarrow$ real container networking.
    \item seccomp, capabilities, AppArmor $\rightarrow$ shrink what the container can do.
    \item \alert{\texttt{-{}-init}} — the PID-1 fix we just met.
  \end{itemize}
\end{frame}

% ============================================================ GOBOXD BRIDGE
\section{The bridge: goboxd}

\begin{frame}{Same problems — we hit them naively}
  \begin{center}
  \begin{tabular}{ll}
    \toprule
    \textbf{Our bash} & \textbf{goboxd spec hole} \\
    \midrule
    \cmd{pivot\_root} + \cmd{umount}   & path traversal \\
    run a process, not \cmd{bash -c "..."} & shell-style commands \\
    user ns + \cmd{-{}-map-root-user}   & UID collisions under load \\
    (not handled)                       & unbounded output — \cmd{rlimit} \\
    \cmd{trap} cleanup                  & stale jail directories \\
    \bottomrule
  \end{tabular}
  \end{center}
\end{frame}

\begin{frame}{nsjail = our script, hardened}
  \begin{itemize}
    \item Everything we wrote, in $\sim$6000 lines of C.
    \item \alert{plus} seccomp, capability dropping, a config language.
    \item \textbf{goboxd}: wrap it in a Go service that runs untrusted code.
  \end{itemize}
\end{frame}

% ============================================================ CLOSE
\begin{frame}{Your turn}
  \begin{itemize}
    \item Catch the \alert{goboxd} finals — tomorrow. You've got the mental model now.
    \item Repo \& slides: \cmd{github.com/sayan-iitm/minidocker}
    \item Stuck / curious: GitHub Discussions on the repo.
  \end{itemize}
\end{frame}

\begin{frame}[standout]
  Questions?
\end{frame}

\end{document}
